\documentclass[a4paper,11pt]{extarticle}
\usepackage[margin=11mm]{geometry}
% Compile with: pdflatex LowLevelFEM_DSL_CheatSheet_pdflatex.tex
% UTF-8 source; DejaVu's Type 1 fonts preserve the original typography.
\usepackage[T1]{fontenc}
\usepackage[utf8]{inputenc}
\usepackage{lmodern} % Scalable math fonts, including small subscripts.
\usepackage[scaled=0.8]{DejaVuSans}
\usepackage[scaled=0.8]{DejaVuSansMono}
\renewcommand{\familydefault}{\sfdefault}
\usepackage{multicol,xcolor,amsmath,amssymb,titlesec,fancyvrb,upquote,graphicx,hyperref}
\input{glyphtounicode}
\pdfgentounicode=1
% These mappings also apply inside the fancyvrb Julia environments.
% Keep the actual Unicode characters in the Julia examples below.
% Match the 0.79 scale of the original monospaced font.
\newcommand{\DSLsymbol}[1]{\scalebox{0.79}{\ensuremath{#1}}}
\DeclareUnicodeCharacter{2022}{\textbullet}
\DeclareUnicodeCharacter{222B}{\DSLsymbol{\int}}
\DeclareUnicodeCharacter{22C5}{\DSLsymbol{\cdot}}
\DeclareUnicodeCharacter{03A9}{\DSLsymbol{\Omega}}
\DeclareUnicodeCharacter{0393}{\DSLsymbol{\Gamma}}
\DeclareUnicodeCharacter{2218}{\DSLsymbol{\circ}}
\DeclareUnicodeCharacter{2207}{\DSLsymbol{\nabla}}
\DeclareUnicodeCharacter{2202}{\DSLsymbol{\partial}}
\DeclareUnicodeCharacter{03BD}{\DSLsymbol{\nu}}
\definecolor{blue}{HTML}{17618A}
\definecolor{gray}{HTML}{52616B}
\hypersetup{colorlinks=true,urlcolor=blue}
\setlength{\columnsep}{7mm}\setlength{\parindent}{0pt}\setlength{\parskip}{2pt}
\titlespacing*{\section}{0pt}{5pt}{2pt}
\titleformat{\section}{\color{blue}\bfseries\fontsize{9}{10}\selectfont}{}{0pt}{}
\newcommand{\code}[1]{\mbox{\ttfamily\small #1}}
\DefineVerbatimEnvironment{julia}{Verbatim}{fontsize=\fontsize{8.6}{8.6}\selectfont,commandchars=\\\{\}}
\begin{document}
{\color{blue}\fontsize{17}{19}\selectfont\bfseries LowLevelFEM.jl}\hfill {\color{gray}\small QUICK REFERENCE • 2026}\par
{\large\bfseries Weak-Form DSL Cheat Sheet}\hfill\small Julia • A4 / one page\par
\vspace{1mm}\hrule\vspace{2mm}
\begin{multicols}{2}
	
	
\paragraph{Workflow} Geometry $\to$ FE mesh (Gmsh) $\to$ fields $\to$ weak form $\to$ sparse matrix, right hand side (rhs) $\to$ boundary conditions (bc) $\to$ solver $\to$ rusulting fields $\to$ post processor $\to$ visualization (Gmsh, Plots).

Field represents both the trying and test space.
\paragraph{Setup} Load LowLeveleFEM package
\begin{julia}
using LowLevelFEM
\end{julia}
%
\paragraph{Geometry} Created by Gmsh or imported STEP or IGES files.
\begin{julia}
openGeometry("part.geo") # opens a *.geo file in Gmsh
openPreProcessor()       # launch Gmsh to draw or make FE mesh
\end{julia}
%
\paragraph{Define the problem} Select the region to solve on, give the type of field to solve
\begin{julia}
mat = Material("part". E=2e5, ν=0.3)
# 3D vector field
U = Field([mat], type=:VectorField, dim=3, fieldName=:u) named "u"
# 2D scalar field
P = Field([mat], type=:ScalarField, dim=2, fieldName=:p) named "p"
\end{julia}
part is a name of a physical group in Gmsh.

type can be :ScalarField, :VectorField, :TensorField.

dim is 1, 2 or 3.

fieldName is arbitrary letter. This will be used durin definition of bc.
%
\paragraph{Weak form} Solve the sparse matrix (system matrix?) and rhs from a weak form.
\begin{julia}
K = ∫(SymGrad(U) ⋅ D(:PlaneStress, mat) ⋅ SymGrad(U), Ω="surf")		# Stiffness matrix in Ω, where dimension of Ω is the same as dimension of U 
\end{julia}
Differential operators in weak form

\begin{tabular}{@{}ll@{}}
\code{U} & identity $u$\\
\code{Grad(U)} & $\nabla u$ (scalar or vector)\\
\code{SymGrad(U)} & engineering strain vector\\
\code{Div(U)} & $\nabla\!\cdot u$\\
\code{Curl(U)} & $\nabla\!\times u$\\
\code{TensorDiv(P)} & divergence of tensor\\
\end{tabular}\par
\subparagraph{2D conventions} 

\code{Grad(U)} gives $[u_{x,x},u_{x,y},u_{y,x},u_{y,y}]$

\code{SymGrad(U)} gives $[u_{x,x},u_{y,y},u_{x,y}+u_{y,x}]$

\subparagraph{3D convention}

\code{SymGrad(U)}: $[xx,yy,zz,xy,yz,zx]$ (engineering shear).
\section{3. Algebra and integration}
\begin{julia}
K = ∫(Grad(P) ⋅ k ⋅ Grad(P); Ω="body")
M = ∫(P ⋅ rho ⋅ P; Ω="body")
f = ∫(P ⋅ q; Ω="body")
g = ∫(P ⋅ flux; Γ="right")
\end{julia}
\code{⋅} contracts operators/coefficient chains; coefficients may be numbers, matrices, or spatially varying fields. \code{Ω} denotes a volume/domain physical group, \code{Γ} a boundary group.
\begin{julia}
K = ∫(Grad(P) ⋅ Grad(P);
      Ω="body", gauss=:full,
      threads=:auto, assembly=:csc)
# Alternative: gauss=:reduced,
# assembly=:ijv, threads=1
\end{julia}
\textbf{Shortcuts:} \code{∫Ω("body", expr)}, \code{∫Γ("right", expr)}. Default matrix assembly is sparse CSC; integration returns \code{SystemMatrix} or an assembled load field.
\section{4. Canonical PDE recipes}
\textbf{Diffusion / Poisson} $-\nabla\cdot(k\nabla T)=q$:
\begin{julia}
K = ∫(Grad(P) ⋅ k ⋅ Grad(P); Ω="body")
f = ∫(P ⋅ q; Ω="body")
\end{julia}
\textbf{2D linear elasticity} $\int\varepsilon(v)^TC\varepsilon(u)$:
\begin{julia}
C = D(:PlaneStress, E, nu)
K = ∫(SymGrad(Pu) ⋅ C ⋅ SymGrad(Pu);
      Ω="body")
f = ∫(Pu ⋅ [0, -1]; Γ="right")
\end{julia}
\columnbreak
\section{5. Solve + essential boundary conditions}
\begin{julia}
bc = BoundaryCondition("left", ux=0, uy=0)
u = solveField(K, f; support=[bc])
# Scalar problem: use a scalar-compatible BC.
\end{julia}
\textbf{Interpretation:} distributed loads and natural boundary terms belong in $f$; Dirichlet constraints are imposed by \code{support}. Other solver options include \code{mpc}, \code{coordSys}, \code{solver}, and \code{preconditioner}.
\section{6. Mixed and coupled fields}
\textbf{Mixed block terms} (e.g. velocity--pressure):
\begin{julia}
# Pu and Pp are distinct Problem descriptors.
B = ∫(Div(Pu) ⋅ Pp; Ω="body")
Bt = ∫(Pp ⋅ Div(Pu); Ω="body")
\end{julia}
\textbf{Compound operator} $B=A_1\nabla u+A_2u$:
\begin{julia}
B = A1 ⋅ Grad(Pu) + A2 ⋅ Pu
K = ∫(B' ⋅ C ⋅ B; Ω="body")
\end{julia}
Compound expressions can generate coupled matrix blocks automatically when operators reference multiple \code{Problem}s. Keep test/trial field ordering consistent.
\section{7. Useful additional forms}
\begin{julia}
# Mass / reaction term
M = ∫(Pu ⋅ rho ⋅ Pu; Ω="body")
# Elastic foundation on a boundary
Kb = ∫(Pu ⋅ [kx 0; 0 ky] ⋅ Pu;
       Γ="edge")
# Divergence penalty
Kd = ∫(Div(Pu) ⋅ Div(Pu); Ω="body")
# Curl-curl form
Kc = ∫(Curl(Pu) ⋅ Curl(Pu); Ω="body")
\end{julia}
\section{8. Field algebra (postprocessing, NOT DSL)}
\begin{julia}
# u and T below are computed FE fields.
eps = (u ∘ ∇ + ∇ ∘ u) / 2
sigma = C * eps  # compatible tensor law
heatflux = -k * ∇(T)
divu = div(u); curlu = curl(u)
\end{julia}
Field operations also include \code{grad}, \code{trace}, \code{∂x}, \code{∂y}, \code{∂z}; the returned fields are distinct from symbolic \code{Grad(P)} / \code{Div(P)} used for matrix assembly.
\section{9. Performance and pitfalls}
\begin{julia}
pat = build_csc_pattern(Pu, Pu; Ω="body")
fill!(pat.nzval, 0.0)
K = ∫(SymGrad(Pu) ⋅ C ⋅ SymGrad(Pu);
      Ω="body", csc_matrix=pat)
\end{julia}
\textbf{Pattern reuse:} clear \code{nzval} before an independent assembly; do \emph{not} call \code{dropzeros!} on the reusable pattern. Copy results before reusing the same matrix storage. Check coefficient/operator dimensions, physical-group names, integration order and shear conventions.
\vspace{2mm}\hrule\vspace{1mm}
{\footnotesize\textbf{Documentation:} \href{https://perebalazs.github.io/LowLevelFEM.jl/dev/reference/multifield/}{DSL / Multifield API} \quad \href{https://perebalazs.github.io/LowLevelFEM.jl/dev/tutorials/}{Tutorials}\\
\textbf{Source:} \href{https://github.com/perebalazs/LowLevelFEM.jl}{github.com/perebalazs/LowLevelFEM.jl}\quad\textit{Examples assume valid mesh, material, and compatible field types.}}
\end{multicols}
\end{document}

